%%%%%%%%%%  scan idx 654  |  folio 641  %%%%%%%%%%
\SGAPrintPage{641}%
of Kummer theory the following fact, due to Grothendieck (SGA 1 XI, on the last page): if \(X\) is a proper scheme over a field \(k\), with \(H^0(X,{O}_X)=k\), then there exists an isomorphism of \(k\)-group schemes \(\underline{\Hom}_{\mathrm{gr}}(\mu_n,\underline{\Pic}_{X/k})\simeq \underline{H}^1(X,\mathbb Z/n\mathbb Z)\); consequently, if \(k\) is algebraically closed, we have \(\Hom_{\mathrm{gr}}(\mu_n,\underline{\Pic}_{X/k})\simeq H^1(X,\mathbb Z/n\mathbb Z)\).
\item If \(k\) has characteristic zero, then the kernel \(N\) is zero, since every finite unipotent group is then zero (SGA 3 XVII 1.7). One thus approximately recovers Kodaira's ``Vanishing theorem'' in the form given by Mumford [2].
\end{enumerate}

\sgaStableTarget{sga6.e13.thm.3-13}\textbf{Theorem 3.13.} \emph{Let \(S\) be a noetherian scheme and \(X\) a projective \(S\)-scheme, endowed with a relatively ample invertible sheaf \(H={O}_X(1)\). We suppose that the geometric fibers of \(X/S\) are integral and of dimension \(r\). Then, for a family \(\mathcal L\) of classes of invertible sheaves on the fibers of \(X/S\), the following conditions are equivalent}~:
\begin{enumerate}[label=(\roman*)]
\item \emph{\(\mathcal L\) is a bounded family.}
\item \emph{In the Hilbert polynomials \(\chi(L_K(n))=\sum_{i=0}^{r}a_i\binom{n+i}{i}\) of the \(L_K\in\mathcal L\), the coefficient \(a_{r-1}\) remains bounded and the coefficient \(a_{r-2}\) remains bounded below.}
\item \emph{As \(L_K\) ranges over \(\mathcal L\), the degree \(d(L_K)=\langle c_1(L_K)\cdot c_1(H_K)^{r-1}\rangle\) remains bounded and the intersection number \(\langle c_1(L_K)^2\cdot c_1(H_K)^{r-2}\rangle\) remains bounded below.}
\item \emph{There exists an integer \(N\) such that we have}
\[
H^{-N}<\mathcal L<H^N
\]
\emph{in the sense of the order relation defined by ampleness\,; that is, for every \(L_K\in\mathcal L\), \(H_K^{\otimes N}\otimes L_K\) and \(H_K^{\otimes N}\otimes L_K^{-1}\) are ample.}
\end{enumerate}

%%%%%%%%%%  scan idx 655  |  folio 642  %%%%%%%%%%
\SGAPrintPage{642}%
Let us prove that (i) implies (ii), (iii), and (iv). Thus suppose that \(\mathcal L\) is bounded by an invertible sheaf on an \(X_T\). Then \(\mathcal L\) satisfies (iv) by EGA II 4.6.12. Suppose \(L\) flat over \(T\). Then the polynomial in \(m\) and \(n\), \(\chi(L_t^{\otimes m}(n))\), depends only on the connected component of \(T\) (EGA III 7.9.4). Consequently, \(\mathcal L\) satisfies (ii) and, taking 2.1 into account, (iii).

Next, note that the family of classes of the \(H_K\) (resp. \({O}_{X_K}\)) is bounded, hence satisfies (iii) (resp. (ii)).

Suppose that \(\mathcal L\) satisfies (iv). Then \(\langle [N\cdot c_1(H_K)+c_1(L_K)]\cdot c_1(H_K)^{r-1}\rangle>0\), hence \(d(L_K)>-Nd(H_K)\), and likewise \(Nd(H_K)>d(L_K)\). Moreover, \(\langle [N c_1(H_K)+c_1(L_K)]^2\cdot c_1(H_K)^{r-2}\rangle>0\), hence \(\langle c_1(L_K)^2\cdot c_1(H_K)^{r-2}\rangle>-2Nd(L_K)-N^2d(H_K)\). Thus \(\mathcal L\) satisfies (iii).

Finally, let us prove that (ii) (resp. (iii)) implies (i). By 2.10 and 2.3, one may replace \(H\) by a multiple\,; therefore suppose \(H\) very ample. Evidently one may decompose \(S\) into pieces\,; thus suppose that there exists a linear section \(Y\) of \(X\) which is integral over \(S\), with fibers of dimension \(2\). By 3.8, one may replace \(X,\mathcal L\) by \(Y,\mathcal L|Y\), (taking account of 1.7 (resp. the projection formula), which implies that condition (ii) (resp. (iii)) on \(X,\mathcal L\) is equivalent to (ii) (resp. (iii)) on \(Y,\mathcal L|Y\))\,; therefore suppose \(r=2\). If \(\delta=\sup\{d(L_K)\}\) is \(>0\), one may replace \(\mathcal L\) by the family of classes of the \(L_K(-\delta)\) by 2.3 and 2.10\,; thus suppose \(\delta\le0\).

If, in the Hilbert polynomials \(\sum_{i=0}^{2}a_i\binom{n+i}{i}\), \(a_1\) remains bounded and \(a_0\) remains bounded below, then the three coefficients remain bounded by 2.3 and 2.7\,; thus \(\mathcal L\) is bounded by 2.11\,; hence (ii) implies
%%%%%%%%%%  scan idx 656  |  folio 643  %%%%%%%%%%
\SGAPrintPage{643}%
(i). On the other hand, (iii) implies (ii) by 2.3 and 2.8.

\sgaStableTarget{sga6.e13.sec.4}\subsection*{4. Finiteness theorems for $\underline{\Pic}^\tau_{X/S}$}

\sgaStableTarget{sga6.e13.para.4-1}\textbf{4.1.} Let \(S\) be a scheme and \(G\) a commutative group functor on \(S\). We denote by \(G^0\) (resp. \(G^\tau\)) the group subfunctor whose \(T\)-points \(f:T\to G\) are those of \(G\) such that, for every algebraically closed field \(K\) and every \(K\)-point \(t\) of \(T\), there exist a connected \(K\)-scheme \(Z\), two \(K\)-points \(z,z_0\) of \(Z\) and a morphism \(g:Z\to G\) such that \(g(z)=f(t)\) and \(g(z_0)=0\) in \(G(K)\) (resp. there exists a nonzero integer \(n\) such that \(n.f(t)\in G^0(K)\)).

It is evident that the formation of \(G^0\) (resp. \(G^\tau\)) commutes with every base change, that a morphism of group functors \(G\to G'\) induces a morphism \(G^0\to G'^0\) (resp. \(G^\tau\to G'^\tau\)), and that if \(S\) is a field and \(G\) is represented by a group scheme locally of finite type, then \(G^0\) is represented by the connected component of the identity element, which is an open (and closed) subscheme, of finite type and geometrically irreducible (SGA 3 VI\(_A\) 2.4).

\sgaStableTarget{sga6.e13.lem.4-2}\textbf{Lemma 4.2.} \emph{Let \(\Phi:G\to G'\) be a morphism between two commutative group functors on a scheme \(S\). We suppose that, for every algebraically closed field \(K\) and every morphism \(\Spec(K)\to S\), \(G_K\) (resp. \(G'_K\)) is representable by a group scheme locally of finite type, and \(\Phi_K\) is representable by a morphism of finite type, (that is, quasi-compact). Then}
\[
  \Phi^{-1}(G'^\tau)=G^\tau.
\]

%%%%%%%%%%  scan idx 657  |  folio 644  %%%%%%%%%%
\SGAPrintPage{644}%
Indeed, we have already noted the inclusion \(\Phi^{-1}(G'^\tau)\supset G^\tau\). For the reverse inclusion, note that \(g\in G(K)\) belongs to \(G^\tau(K)\) if and only if the sequence of \(ng\) (\(n\in\mathbb N\)) is quasi-compact\,; since \(\Phi_K\) is quasi-compact, it suffices for this that the sequence of \(\Phi(ng)=n\Phi(g)\) have the same property\,; this gives the assertion.

\sgaStableTarget{sga6.e13.lem.4-3}\textbf{Lemma 4.3.} \emph{Let \(X\), \(Y\) be proper over a scheme \(S\) and \(f:X\to Y\) a surjective morphism (resp. the inclusion morphism of the zero scheme of a section \(\varphi\) of an ample sheaf \({O}_Y(1)\), under the hypotheses of 3.8). Then}
\[
  (f^*)^{-1}\underline{\Pic}^\tau_{X/S}=\underline{\Pic}^\tau_{Y/S}.
\]

Indeed, the assertion is a formal consequence of 4.2, (XII 1.2), 3.1, and 3.5 (resp. 3.8).

\sgaStableTarget{sga6.e13.para.4-4}\textbf{4.4.} Let \(k\) be an algebraically closed field, \(X\) a proper \(k\)-scheme, and \(L\) an invertible sheaf on \(X\). We say that \(L\) is \uline{algebraically equivalent} to \(0\) (resp. \(\tau\)-\uline{equivalent} to \(0\)) if there exist a connected \(k\)-scheme \(T\) of finite type, an invertible sheaf \(M\) on \(X_T=X\times_k T\), and two \(k\)-points \(t,t_0\) of \(T\) such that \(L\simeq M_t\) and \({O}_X\simeq M_{t_0}\) (resp. if there exists a nonzero integer \(n\) such that \(L^{\otimes n}\) is algebraically equivalent to \(0\)), or, what is easily seen to be equivalent, if the \(k\)-point \(\lambda\) of \(\underline{\Pic}_{X/k}\) corresponding to \(L\) belongs to \(\underline{\Pic}^0_{X/k}(k)\) (resp. to \(\underline{\Pic}^\tau_{X/k}(k)\)).

Recall also that \(L\) is said to be \uline{numerically equivalent} to \(0\) if, for every integral closed curve \(Y\) in \(X\), the intersection number \(\langle c_1(L)\cdot Y\rangle\) is \(0\) (see 2.1 or X 4.5).

%%%%%%%%%%  scan idx 658  |  folio 645  %%%%%%%%%%
\SGAPrintPage{645}%
\sgaStableTarget{sga6.e13.prop.4-5}\textbf{Proposition 4.5. (Functorial behavior).} \emph{Let \(f:X'\to X\) be a morphism between two proper \(k\)-schemes, and let \(L\) be an invertible module on \(X\). Then}~:
\begin{enumerate}[label=(\roman*)]
\item \emph{If \(L\) is \(\tau\)-equivalent to \(0\) (resp. numerically equivalent to \(0\)), then so is \(f^*L\).}
\item \emph{Conversely, when \(f\) is surjective, if \(f^*L\) is \(\tau\)-equivalent to \(0\) (resp. numerically equivalent to \(0\)), then so is \(L\).}
\end{enumerate}

Indeed, assertion (i) (resp. (ii)) for \(\tau\)-equivalence follows at once from 4.1 (resp. 4.3). On the other hand, let \(Y'\) be an integral closed curve in \(X'\), and \(Y=f(Y')\) its image, endowed with the induced reduced structure. When \(\dim(Y)=1\), we have the projection formula (X 4.4.4)
\sgaStableTarget{sga6.e13.eq.4-5-1}\[
  \tag*{4.5.1.}
  \langle c_1(f^*L)\cdot Y'\rangle
  =[k(Y'):k(Y)]\,\langle c_1(L)\cdot Y\rangle
\]
where \(k(Y')\) (resp. \(k(Y)\)) is the function field of \(Y'\) (resp. \(Y\)), and when \(\dim(Y)=0\), we have \((f^*L)_{Y'}\simeq{O}_{Y'}\). Consequently, if \(L\) is numerically equivalent to \(0\), then \(\langle c_1(f^*L)\cdot Y'\rangle=0\)\,; hence (i).

Conversely, when \(f\) is surjective, given an integral closed curve \(Y\) in \(X\), one can find an integral closed curve \(Y'\) in \(X'\) such that \(Y=f(Y')\): let \(y\) be the generic point of \(Y\), \(y'\) a point of the fiber \(f^{-1}(y)\) algebraic over \(k(y)\), and \(Y'\) the closure of \(y'\). Then 4.5.1 shows that \(\langle c_1(f^*L)\cdot Y'\rangle=0\) implies \(\langle c_1(L)\cdot Y\rangle=0\)\,; hence (ii).

\sgaStableTarget{sga6.e13.thm.4-6}\textbf{Theorem 4.6. (The numerical characterization of \(\underline{\Pic}^\tau_{X/k}\)).} \emph{Let \(k\) be an algebraically closed field, \(X\) a proper \(k\)-scheme, and \(L\) an invertible sheaf on \(X\). Then the following conditions are equivalent}~:
\begin{enumerate}[label=(\roman*)]
\item \emph{\(L\) is \(\tau\)-equivalent to \(0\).}
\item \emph{For every coherent sheaf \(F\) on \(X\), we have}
